\begin{center}
  \textsc{Sommario}
\end{center}
%
\noindent
%
La decomposizione CP (Canonical Polyadic Decomposition) di rango $ r $ di un tensore è la sua
scomposizione nella somma di $ r $ tensori di rango uno, una generalizzazione naturale della
decomposizione di matrici. Oltre ad avere un'importanza teorica nel calcolo tensoriale, è uno strumento utile in
molte aree scientifiche: per citarne alcune, spettroscopia, \emph{signal processing}, neurologia e
analisi dei dati. In
molte di queste applicazioni l'unicità della fattorizzazione è una proprietà fondamentale per rappresentare
fedelmente le componenti di sistemi complessi---si pensi al problema di modellare i 
neuroni di un sistema nervoso tramite un tensore per studiarne gli impulsi, noto come \emph{blind
source separation}. 

Il primo obiettivo è dimostrare l'unicità della fattorizzazione CP sotto
opportune ipotesi.
Questo risultato, originariamente dimostrato da Kruskal nel 1977 attraverso il concetto di Kruskal
rank, è stato di recente semplificato seguendo la dimostrazione originale, usando strumenti
geometrici. Inoltre, osserveremo alcune peculiarità teoriche che
distinguono i tensori dalle matrici, ponendo le basi per l'utilizzo del calcolo tensoriale nella
teoria della complessità.

Successivamente, applicheremo questi strumenti al problema della complessità della moltiplicazione
tra matrici: in quanto mappa bilineare, il prodotto matriciale può essere rappresentato tramite un
tensore tre indici, e poi messo in corrispondenza con un algoritmo per valutare il prodotto in
funzione delle entrate del tensore. Introdurremo la nozione di complessità aritmetica del prodotto
tra matrici per definire
l'esponente asintotico $ \omega $, che ne misura la complessità ottimale. Tale quantità sarà messa in corrispondenza con il rango CP
di un tensore tramite il Teorema di Strassen. Come corollario, ridurre lo studio della complessità ai tensori, anziché all'analisi
diretta delle equazioni algebriche, semplifica notevolmente la ricerca di algoritmi più efficienti;
una sfida aperta da decenni, dall'algoritmo di Strassen (del 1970), fino alle recenti scoperte
guidate dall'intelligenza artificiale. 

Infine, l'analisi verrà estesa dagli algoritmi di moltiplicazione esatti agli algoritmi approssimati
(APA), introdotti da Bini et al. In questo contesto emerge il fenomeno topologico del border rank,
per il quale una successione di tensori di rango inferiore può convergere a un tensore di rango
superiore. Dal punto di vista algoritmico, ciò si traduce nella possibilità di approssimare la
moltiplicazione esatta tramite una successione di algoritmi con
costo computazionale inferiore. Di conseguenza,  il
rango tensoriale può essere sostituito con il border rank nel calcolo di $ \omega $, permettendo di ottenere stime asintotiche più
accurate, al prezzo di un errore algoritmico arbitrariamente piccolo.
